% Generated draft: requires mathematical and source review.
\documentclass[11pt]{article}
\usepackage[margin=25mm]{geometry}
\usepackage{fontspec}
\setmainfont{Latin Modern Roman}
\usepackage{amsmath,amssymb,mathtools,unicode-math}
\setmathfont{Latin Modern Math}
\usepackage{enumitem,needspace,xurl,hyperref}
\setlength{\parindent}{0pt}
\setlength{\parskip}{5pt}
\setlength{\emergencystretch}{3em}
\newcommand{\N}{\mathbb{N}}
\newcommand{\Z}{\mathbb{Z}}
\newcommand{\Q}{\mathbb{Q}}
\newcommand{\R}{\mathbb{R}}
\newcommand{\C}{\mathbb{C}}
\newcommand{\F}{\mathbb{F}}
\newcommand{\A}{\mathbb{A}}
\newcommand{\PP}{\mathbb{P}}
\newcommand{\E}{\mathbb{E}}
\newcommand{\eps}{\varepsilon}
\newcommand{\dd}{\,\mathrm{d}}
\begin{document}
\section*{UnsolvedMath ID: 30004875}
\textbf{Source problem code:} catalogue entry
\subsection{Definitions and mathematical statement}
For a field \(k\) and a \(k\)-algebra \(A\), a finite-dimensional representation means a left \(A\)-module \(M\) with \(\dim_k M<\infty\). Such a module is indecomposable if \(M\ne 0\) and every decomposition \(M\cong M_1\oplus M_2\) into \(A\)-modules has \(M_1=0\) or \(M_2=0\). These definitions explain the terminology; the supplied record does not specify the field or the quantum loop algebras intended.

The supplied statement is: ``Classify the finite-dimensional indecomposable representations of quantum loop algebras; equivalently, split their Grothendieck groups.'' Under the conventional interpretation that classification is up to module isomorphism, its recoverable classification target, for each intended algebra \(A\) over its intended field \(k\), is to describe
\[
\left\{[M]_{\cong}: M\text{ is an indecomposable left }A\text{-module},\ \dim_k M<\infty\right\}.
\]
The field, underlying Lie-theoretic data, deformation parameter and its restrictions, version of the quantum loop algebra, and permitted representation category are unspecified. Thus this is a conditional description of the target, not a fully determined problem with supplied hypotheses.

For an additive category \(\mathcal C\) of finite-dimensional modules, its split Grothendieck group is the abelian group generated by isomorphism classes \([M]\), subject to \([M\oplus N]=[M]+[N]\). The ordinary Grothendieck group instead imposes \([M]=[L]+[N]\) for short exact sequences \(0\to L\to M\to N\to0\). The phrase ``split their Grothendieck groups'' does not specify which group is meant or what splitting is requested. Its asserted equivalence with classifying indecomposables cannot be adopted without clarification.

The retained original passage reports a discussion leading to an ``Open problem'' but contains no mathematical question. Although the registry labels the replacement statement \texttt{corrected\_verified} and reports verification against the official report, the supplied excerpt does not establish its hypotheses or the asserted equivalence. These source limitations remain unresolved.
\subsection{Short English statement}
The supplied target is to classify finite-dimensional indecomposable representations of quantum loop algebras, conventionally up to isomorphism. The algebraic parameters and representation category are missing, and the claimed equivalent task concerning Grothendieck groups is unclear. The supplied open label and verification claims require human review against the report.
\subsection{Sources}
\begin{itemize}
\item Supplied local UnsolvedMath registry record, permanent ID \(30004875\), title \emph{Indecomposable Representations of Quantum Loop Algebras}, exact source problem code \texttt{catalogue entry}; collection: \emph{Oberwolfach Reports Open Problems}, category: Algebra.
\item Supplied source citation: \emph{Mini-Workshop: Three Facets of R-Matrices (2022)}, Oberwolfach Reports, pp.~2791--2825, DOI \texttt{10.4171/owr/2021/51}. Source URL and supplied literature reference labelled ``OWR report'': \url{https://doi.org/10.4171/owr/2021/51}. This authoritative report is cited by the record; its contents have not been independently consulted here.
\item Supplied attribution: the retained passage says that discussion between David Hernandez and Valerio Toledano Laredo (participating online) led to the open problem. No individual proposer is specified. The proposed year and year tag are 2021, while the supplied report citation displays 2022; this dating distinction is retained without resolution.
\item Literature assessment checked at \texttt{2026-08-22}: supplied dated curation, not independently verified. It states that the extracted statement contains only the heading ``Open problem'' and omits the mathematical question, making resolution indeterminate and leading to classification as open. The background expressly requests independent verification of time-sensitive status and cited claims.
\item The record separately supplies status \texttt{open}, statement status \texttt{corrected\_verified}, and a claim that the clean statement was checked against the official report and preserves its hypotheses, notation, and target. These are registry assertions, not independently verified mathematical facts.
\end{itemize}
\end{document}
